%-------------------------------------------------------------------------------
\section{Spectral analysis tool}
\label{sec:postprocess_spectral_analysis}
%-------------------------------------------------------------------------------

\verb|spectral_analysis| is a post-processing tool for calculating horizontal spectra from the output files of \scaledg.
It reads variables in the SCALE-netCDF files and outputs spectral coefficients and one-dimensional spectra to a NetCDF file.
This tool is useful, for example, for analyzing kinetic energy spectra at specified vertical levels.
Currently, this tool is available only for the output data of regional simulations on the Cartesian coordinates.

To compile \verb|spectral_analysis|, execute the following command under the top directory after the compilation of \scaledg:
%%
\begin{verbatim}
$ cd model/atm_nonhydro3d/util/spectral_analysis/
$ make
\end{verbatim}
%%
If the compilation succeeds, the executable binary file is generated under the \verb|bin| directory.

The basic settings are specified in \namelist{PARAM_SPECTRAL_ANALYSIS}.
For example,

\editboxtwo{
\verb|&PARAM_SPECTRAL_ANALYSIS | & \\
\verb|in_bs_filebase      = "init_00010101-000000.000",| & ; Base name of input file for basic state \\
\verb|in_filebase         = "history",|                  & ; Base name of input history files \\
\verb|out_filebase        = "energy_spectra",|           & ; Base name of output file \\
\verb|target_proc_num_tot = 32,|                         & ; Total number of MPI processes used in the original run \\
\verb|LevelNum            = 1,|                          & ; Number of target vertical levels \\
\verb|TARGET_LEVELS       = 800.D0,|                     & ; Target vertical levels \\
\verb|level_units         = "m",|                        & ; Unit of target vertical levels \\
\verb|VARS                = "U", "V", "W",|              & ; Target 3D variables \\
\verb|SpectralAnalysisType = "Uniform_Sampling",|        & ; Type of spectral analysis \\
\verb|ST_NsamplePtPerElem1D = 8,|                        & ; Number of sampling points per element \\
\verb|ST_ks = -64, ST_ke = 64,|                          & ; Range of wavenumber in the first direction \\
\verb|ST_ls = -64, ST_le = 64,|                          & ; Range of wavenumber in the second direction \\
\verb|KinEnergyAnalysisFlag = .true.,|                   & ; Output kinetic energy spectra \\
\verb|/ | & \\
}
%%
\nmitem{in_filebase} specifies the base name of the input history files, and \nmitem{out_filebase} specifies the base name of the output NetCDF file.
\nmitem{target_proc_num_tot} should be set to the total number of MPI processes used in the original simulation.
\nmitem{TARGET_LEVELS} and \nmitem{level_units} specify the vertical levels at which three-dimensional variables are analyzed.
\nmitem{VARS} specifies the target three-dimensional variables.
Two-dimensional variables can also be analyzed by specifying \nmitem{VARS2D}.
\nmitem{SpectralAnalysisType} can be set to \verb|Uniform_Sampling| or \verb|L2_Projection|.


The mesh information of the original output is specified in
\namelist{PARAM_SPECTRAL_ANALYSIS_MESH}.

\editboxtwo{
\verb|&PARAM_SPECTRAL_ANALYSIS_MESH | & \\
\verb|dom_xmin = -1.6D3, dom_xmax = 1.6D3,| & ; Domain range in the $x$-direction \\
\verb|dom_ymin = -1.6D3, dom_ymax = 1.6D3,| & ; Domain range in the $y$-direction \\
\verb|dom_zmin = 0.D0,   dom_zmax = 1.6D3,| & ; Domain range in the $z$-direction \\
\verb|NprcX = 8,|                            & ; Number of MPI processes in the $x$-direction in the original run \\
\verb|NeX   = 2,|                            & ; Number of elements in the $x$-direction in each MPI process \\
\verb|NprcY = 4,|                            & ; Number of MPI processes in the $y$-direction in the original run \\
\verb|NeY   = 4,|                            & ; Number of elements in the $y$-direction in each MPI process \\
\verb|NeZ   = 8,|                            & ; Number of elements in the $z$-direction \\
\verb|PolyOrder_h = 7,|                      & ; Polynomial order in the horizontal direction \\
\verb|PolyOrder_v = 7,|                      & ; Polynomial order in the vertical direction \\
\verb|/ | & \\
}


The execution of \verb|spectral_analysis| is as follows:
%%
\begin{verbatim}
$ mpirun -n 4 ./spectral_analysis spectral_analysis.conf
\end{verbatim}
%%
In this example, \verb|spectral_analysis| is executed with four MPI processes. 
The last argument specifies the configuration file.
Note that the number of MPI processes used is constrained by the current implementation.
\begin{itemize}
    \item The number of MPI processes used for \verb|spectral_analysis| must be a divisor of \nmitem{target_proc_num_tot} and must not exceed \nmitem{target_proc_num_tot}.
    \item When \nmitem{SpectralAnalysisType} is set to \verb|Uniform_Sampling|, the number of MPI processes used for \verb|spectral_analysis| should be the same as \nmitem{NprcY}.
\end{itemize}


The output file is generated as \verb|[out_filebase].nc|.
The output contains wavenumber coordinates and spectra of the specified
variables. When \nmitem{KinEnergyAnalysisFlag} is set to \verb|.true.|,
horizontal and vertical kinetic energy spectra are also output.